\documentclass[11pt]{article}
\usepackage{mathblatt}
\begin{document}
\weit
\typeout{PROBE-BEGIN scharen-von-geraden-und-ebenen-e1-k2-s6-v1}
\begin{aufgabe}{\detokenize{scharen-von-geraden-und-ebenen-e1-k2-s6-v1}}
\begin{teile}
\teil Probe

\begin{ksys3}\rpunkt{0}{3}{0}{P}\rpunkt{0}{0}{1}{Q}\rgerade[-0.3:1.5]{0,3,0}{0,-3,1}{h}\end{ksys3}
\end{teile}
\end{aufgabe}
\typeout{PROBE-END scharen-von-geraden-und-ebenen-e1-k2-s6-v1}
\typeout{PROBE-BEGIN schnittmengen-e3-k1-s5-v1}
\begin{aufgabe}{\detokenize{schnittmengen-e3-k1-s5-v1}}
\begin{teile}
\teil Probe

\begin{ksys3}[xyz,x1min=0,x1max=8,x2min=0,x2max=8,x3min=0,x3max=5] \rquader{0,0,0}{7}{7}{2} \rgerade[0:1]{7,0,0}{0,3.5,2}{} \rgerade[0:1]{7,3.5,2}{-3.5,3.5,0}{} \rgerade[0:1]{3.5,7,2}{-3.5,0,-2}{} \rgerade[0:1]{0,7,0}{7,-7,0}{} \rgerade[0:5]{7,7,0}{0,0,1}{h} \end{ksys3}
\end{teile}
\end{aufgabe}
\typeout{PROBE-END schnittmengen-e3-k1-s5-v1}
\typeout{PROBE-BEGIN vektoren-und-rechenoperationen-e1-k2-s5-v1}
\begin{aufgabe}{\detokenize{vektoren-und-rechenoperationen-e1-k2-s5-v1}}
\begin{teile}
\teil Probe

\begin{ksys3} \rgerade[0:1]{0,0,0}{4,4,4}{} \rgerade[0:1]{0,4,0}{4,-4,4}{} \end{ksys3} \begin{ksys}[xmin=0,xmax=5,ymin=0,ymax=5,xlabel=$x_1$,ylabel=$x_2$] \end{ksys}
\end{teile}
\end{aufgabe}
\typeout{PROBE-END vektoren-und-rechenoperationen-e1-k2-s5-v1}
\typeout{PROBE-BEGIN scharen-von-geraden-und-ebenen-e1-k2-s6-v2}
\begin{aufgabe}{\detokenize{scharen-von-geraden-und-ebenen-e1-k2-s6-v2}}
\begin{teile}
\teil Probe

\begin{ksys3}\rpunkt{2}{0}{1}{R}\rpunkt{1}{0}{3}{S}\rgerade[-0.5:1.8]{2,0,1}{-1,0,2}{h}\end{ksys3}
\end{teile}
\end{aufgabe}
\typeout{PROBE-END scharen-von-geraden-und-ebenen-e1-k2-s6-v2}
\typeout{PROBE-BEGIN schnittmengen-e3-k1-s5-v2}
\begin{aufgabe}{\detokenize{schnittmengen-e3-k1-s5-v2}}
\begin{teile}
\teil Probe

\begin{ksys3}[xyz,x1min=0,x1max=6,x2min=0,x2max=6,x3min=0,x3max=6] \rquader{0,0,0}{5}{5}{1} \rgerade[0:1]{5,0,0}{0,1,1}{} \rgerade[0:1]{5,1,1}{-4,4,0}{} \rgerade[0:1]{1,5,1}{-1,0,-1}{} \rgerade[0:1]{0,5,0}{5,-5,0}{} \rgerade[0:6]{5,5,0}{0,0,1}{h} \end{ksys3}
\end{teile}
\end{aufgabe}
\typeout{PROBE-END schnittmengen-e3-k1-s5-v2}
\end{document}
